\section{Conclusiones}

La situación actual del planeamiento urbanístico de Santa Cruz de Tenerife se encuentra jurídicamente fragmentada, dado que diversas resoluciones judiciales han anulado sistemáticamente los planes generales más recientes, así como sus adaptaciones provisionales. El Plan General de Ordenación (PGO) fue anulado por el Tribunal Superior de Justicia de Canarias (TSJC) mediante la sentencia 116/2014 \cite{tsjc_116_2014}, una anulación que se publicó oficialmente en octubre de 2020. Tras esta invalidación, el municipio retomó el marco normativo anterior, el Plan General de Ordenación Urbana de 2005 (PGOU-05) \cite{pgou_santacruz_2005}, e intentó actualizarlo mediante la aprobación, en mayo de 2022, de una Ordenanza Municipal Provisional relativa al régimen de usos. No obstante, el TSJC suspendió posteriormente dicha ordenanza provisional de 2022, a través de la sentencia 000070/2023 \cite{tsjc_070_2023}.

Esta inestabilidad en la planificación urbana genera una importante inseguridad jurídica, especialmente en lo que respecta a las licencias de edificación. Sin un plan integral vigente y actualizado, el marco legal para la concesión de licencias se vuelve ambiguo. Si bien algunas fuentes sugieren la aplicación parcial del marco normativo de 2005 o de regulaciones transitorias, la anulación de la ordenanza de 2022 ha socavado, en particular, el intento más reciente de establecer normas claras para el uso del suelo. Esta situación exige extremar las precauciones ante cualquier proyecto urbanístico, ya que la validez de los permisos podría ser impugnada debido a la contradicción entre los estatus jurídicos de los distintos instrumentos de planificación.
